\uuid{pt2w}
\titre{Composée/Quotient/DL usuels}
\auteur{hanani ; mzari}
\organisation{exo7}
\licence{CC BY-NC-SA 4.0}
\datecreate{2020-01-20}
\niveau{L1}
\module{Analyse}
\chapitre{Développement limité}
\sousChapitre{Calculs}
\difficulte{3}
\theme{Opérations sur les DL, Composée, Quotient, DL usuels}
\exo7qcmid{lille-2.developpements-limites.016}

\contenu{
\texte{
Soit $a\in \Rr$ et $f$ une fonction. On dit que $f(x)$ est un $o\left((x-a)^n\right)$ au voisinage de $a$ et on écrit $f(x)=o((x-a)^n)$ si $\displaystyle \lim _{x\to a}\frac{f(x)}{(x-a)^n}=0$. On note par $DL_n(a)f(x)$ le développement limité de $f$ à l'ordre $n$ en $a$. On note aussi $DL_n(a^+)f(x)$ (resp. $DL_n(a^-)f(x)$) le développement limité de $f$ à l'ordre $n$ à droite en $a$ (resp. à gauche en $a$).
}

\begin{question}
\qclassification{125.02}
\qtags{motcle=Composée/Quotient/DL usuels}
Soit $\displaystyle f(x)=\frac{\ln (1+x+x^2)}{\sqrt{1+2x}-1}$. Parmi les affirmations suivantes, cocher celles qui sont vraies :
\begin{answers}
\good{$\displaystyle DL_2(0)\ln (1+x+x^2)=x+\frac{x^2}{2}+o(x^2)$.}
\good{$\displaystyle DL_2(0)\left(\sqrt{1+2x}-1\right)=x-\frac{x^2}{2}+o(x^2)$.}
\bad{$\displaystyle DL_2(0)f(x)=1+x+\frac{x^2}{2}+o(x^2)$.}
\bad{$\displaystyle DL_3(0)[xf(x)]=x+x^2+\frac{x^3}{2}+o(x^3)$.}
\end{answers}
\begin{explanations}
Pour écrire le $DL_2(0)f(x)$, il faut écrire les développements limités du numérateur et du dénominateur à l'ordre $3$ en $0$. Ensuite effectuer la division suivant les puissances croissantes à l'ordre $2$. On aura :
$$DL_2(0)f(x)=1+x-\frac{2x^2}{3}+o(x^2).$$
En multipliant par $x$, on obtient : $\displaystyle DL_3(0)xf(x)=x+x^2-\frac{2x^3}{3}+o(x^3)$.
\end{explanations}
\end{question}
}
